%% =====================================================================
%%  情報処理学会論文誌ジャーナル 投稿用テンプレート（和文・submit）
%% ---------------------------------------------------------------------
%%  コンパイラ : uplatex（または platex）+ dvipdfmx
%%               Overleaf では Menu > Compiler を「LaTeX」、
%%               TeX Live を最新、Main document を main.tex に設定。
%%               （latexmk 用の設定は latexmkrc を参照）
%%  必要ファイル: ipsj.cls, ipsjsort.bst を本プロジェクトに同梱すること。
%%               情報処理学会の「LaTeX スタイルファイル」配布ページから
%%               入手し、このフォルダにアップロードする。
%%  参考文献   : references.bib + ipsjsort.bst（BibTeX）
%% =====================================================================
%%
%% 投稿時は [submit] を付ける。カメラレディ時は [submit] を外す。
\documentclass[submit,techrep]{ipsj}

\usepackage[dvipdfmx]{graphicx}
\usepackage{latexsym}

%% 受付・採録日（投稿時は仮の日付でよい）
\受付{2026}{6}{16}
\採録{2026}{6}{16}

\begin{document}

%% ---- タイトル -------------------------------------------------------
\title{ここに和文タイトルを記入する}
\etitle{English Title Here}

%% ---- 所属 -----------------------------------------------------------
%% \affiliate{ラベル}{和文所属\\英文所属}
\affiliate{IPSJ}{所属機関名\\
Affiliation, City, Country}
%% 現在の所属など（任意）
%\paffiliate{JU}{現在，○○大学}

%% ---- 著者 -----------------------------------------------------------
%% \author{和名}{英名}{所属ラベル}[メールアドレス]
\author{情報 太郎}{Taro Joho}{IPSJ}[author@example.ac.jp]
\author{処理 花子}{Hanako Shori}{IPSJ}

%% ---- 概要・キーワード ----------------------------------------------
\begin{abstract}
ここに和文の概要を記入する．本論文では……について述べ，
……を提案し，……により有効性を示す．
\end{abstract}

\begin{jkeyword}
キーワード1，キーワード2，キーワード3
\end{jkeyword}

\begin{eabstract}
Write the English abstract here. This paper proposes ...
and demonstrates its effectiveness through ...
\end{eabstract}

\begin{ekeyword}
keyword1, keyword2, keyword3
\end{ekeyword}

\maketitle

%% =====================================================================
\section{はじめに}
本文をここに記述する．参考文献は\cite{ipsj-sample}のように引用する．
複数引用は\cite{ipsj-sample,book-sample}のようにまとめられる．

\section{関連研究}

\section{提案手法}
図は次のように挿入する（図~\ref{fig:example}）．

\begin{figure}[tb]
  \centering
  %\includegraphics[width=\columnwidth]{figure.pdf}
  \caption{図のキャプション}
  \label{fig:example}
\end{figure}

表は次のように記述する（表~\ref{tab:example}）．

\begin{table}[tb]
  \caption{表のキャプション}
  \label{tab:example}
  \centering
  \begin{tabular}{lcr}
    \hline
    項目 & 値A & 値B \\
    \hline
    行1 & 1 & 2 \\
    行2 & 3 & 4 \\
    \hline
  \end{tabular}
\end{table}

\section{評価}

\section{おわりに}

%% ---- 謝辞（任意）---------------------------------------------------
%\begin{acknowledgment}
%本研究は……の支援を受けた．
%\end{acknowledgment}

%% ---- 参考文献（BibTeX）--------------------------------------------
\bibliographystyle{ipsjsort}
\bibliography{references}

%% ---- 著者紹介（投稿時は任意・カメラレディ時に記入）---------------
%\begin{biography}
%\profile{m}{情報 太郎}{19XX年生まれ．……}
%\profile{n}{処理 花子}{19XX年生まれ．……}
%\end{biography}

\end{document}
